\documentclass[11pt,a4paper]{article}

\usepackage[utf8]{inputenc}
\usepackage[T1]{fontenc}
\usepackage[french]{babel}
\usepackage{hyperref}
\usepackage{geometry}
\usepackage{xcolor}
\usepackage{enumitem}
\usepackage{titlesec}

% Marges
\geometry{
    top=2.5cm,
    bottom=2.5cm,
    left=2.5cm,
    right=2.5cm
}

% Couleurs
\definecolor{primary}{HTML}{7C4DFF}
\definecolor{secondary}{HTML}{333333}

% Personnalisation des titres
\titleformat{\section}{\color{primary}\normalfont\Large\bfseries}{\thesection.}{1em}{}
\titleformat{\subsection}{\color{secondary}\normalfont\large\bfseries}{\thesubsection.}{1em}{}

% Méta-données
\title{\textbf{\textcolor{primary}{Rapport de Projet : TravelingApp}} \\ \large Bilan final de la fusion et des développements}
\author{Reda \& Binôme}
\date{\today}

\begin{document}

\maketitle
\thispagestyle{empty}

\vspace{1cm}

\section{Présentation du Dépôt et Lien}
Ce dépôt regroupe l'application finale issue de la fusion complète des deux modules principaux développés au cours du semestre :
\begin{itemize}
    \item \textbf{TravelPath (Java) :} Module de planification intelligente de trajets interactifs.
    \item \textbf{TravelShare (Kotlin) :} Module de réseau social dédié au partage d'expériences de voyages.
\end{itemize}

\vspace{0.5cm}
\noindent \textbf{Lien du dépôt Git :} \url{https://github.com/redah22/TravelingApp}
\vspace{1cm}

\section{Parties Développées (Acquises et Fonctionnelles)}

\subsection{Module Social (TravelShare - Partie 1)}
\begin{itemize}[label=\textcolor{primary}{\textbullet}]
    \item \textbf{Persistance via Base de Données SQLite :} Intégration locale robuste (\texttt{DatabaseHelper.kt}) sauvegardant l'intégralité des publications, les images (URI physiques locales Android), les commentaires, et l'inscription sécurisée des utilisateurs.
    \item \textbf{Flux d'actualité :} Création dynamique d'objets \texttt{PhotoPost} interfacés avec des éléments de cartes (Lieu, Auteur, Description).
    \item \textbf{Interactions communautaires :} Appui sur la logique Many-To-Many en SQLite pour autoriser les favoris (Likes dynamiques), ainsi que la publication et consultation dynamique de commentaires avec affichage des compteurs correspondants.
    \item \textbf{Authentification fluide :} Flux UI harmonisé couvrant le \textit{Register}, le \textit{Login} et un accès en mode \textit{Anonyme}.
\end{itemize}

\subsection{Module Planification (TravelPath - Partie 2)}
\begin{itemize}[label=\textcolor{primary}{\textbullet}]
    \item \textbf{Moteur de recherche temps-réel :} Saisie avec autocomplétion asynchrone pointant vers l'API \url{photon.komoot.io} (OpenStreetMap) pour le géocodage des destinations.
    \item \textbf{Générateur d'Aventures asynchrone :} Algorithme multi-paramètres filtrant sur le \textit{budget}, \textit{la durée}, et la \textit{densité d'effort}. L'exploration directe de l'API Overpass génère des Points d'Intérêt (POI) qualitatifs autour de la zone et du secteur d'activité visés.
    \item \textbf{Carte Dynamique et Tuiles OSMDroid :} Rendu visuel d'OpenStreetMap (Mapnik) affichant un tracé (\textit{Polyline}) du parcours suggéré, superposé par des bulles (marqueurs configurés et numérotés) détaillant les distances.
    \item \textbf{Métriques additionnelles :} Connexion dynamique à l'API \textit{Open-Meteo} pour l'affichage précis des prévisions actuelles de la coordonnée GPS cible. Algorithme de \textit{Haversine} permettant de formuler le temps réel de marche entre deux monuments consécutifs.
    \item \textbf{Export document :} Architecture de génération de documents PDF natifs reprenant le planning de la journée heure par heure.
\end{itemize}

\subsection{Logique de Fusion et d'Architecture globale}
\begin{itemize}[label=\textcolor{primary}{\textbullet}]
    \item \textbf{Synchronisation de Sessions asymétriques :} Couplage de la base de données \textit{SQLite} à base de requêtes brutes de TravelShare avec la gestion \textit{SharedPreferences} rapide de TravelPath, assurant à l'utilisateur de n'avoir à créer son compte (\texttt{RegisterActivity}) ou se connecter qu'une seule et unique fois quel que soit l'écran sur lequel l'app démarre.
    \item \textbf{Intégrations Backstack :} Modification du cycle de vie des activités (Kotlin) et des fragments (Java) pour conserver une cohérence de la mémoire. Nettoyage absolu lors du bouton de retour au point central \texttt{MainActivity}.
    \item \textbf{UX/UI harmonisée :} Tous les Layouts Kotlin (Sombre) ont basculé sur un thème \textit{Light Mode DayNight} pour concorder chromatiquement avec l'interface originelle Violet et Blanc (\texttt{\#7C4DFF}) du binôme Java, aboutissant à une charte graphique globale propre et continue.
\end{itemize}

\newpage

\section{Développements de la phase finale}

\subsection{Corrections fonctionnelles (Back-end)}
\begin{itemize}[label=\textcolor{primary}{\textbullet}]
    \item \textbf{Authentification unifiée :} Le \texttt{LoginFragment} (Java) interroge désormais en cascade les deux sources de données : d'abord les \textit{SharedPreferences} (chemin rapide), puis \texttt{DatabaseHelper.kt} via un fallback SQLite si le compte n'est connu que de la base de données (cas d'une réinstallation). La méthode \texttt{getInstance} a été annotée \texttt{@JvmStatic} pour permettre son appel direct depuis Java.
    
    \item \textbf{Filtre de visibilité du fil public :} Ajout de la méthode \texttt{recupererPostsPublics()} dans \texttt{DatabaseHelper} qui ne retourne que les publications ayant le scope \texttt{PUBLIC}. Le fil d'actualité (\texttt{FeedActivity}) et la carte (\texttt{MapFeedActivity}) utilisent désormais cette méthode, garantissant que les posts privés et de groupe n'apparaissent pas dans l'espace public.
    
    \item \textbf{Contexte utilisateur dans la carte :} \texttt{MapFeedActivity} transmettait systématiquement \texttt{IS\_ANONYMOUS = true} aux fiches détaillées. Ce bug empêchait les utilisateurs connectés d'interagir depuis la carte. Le contexte est maintenant passé correctement via les extras de l'Intent.
    
    \item \textbf{Synchronisation du fil après publication :} Correction d'un bug de synchronisation entre \texttt{postsList} et \texttt{currentDisplayList} dans \texttt{FeedActivity} : les nouvelles publications sont maintenant insérées dans les deux listes, assurant la cohérence de la recherche immédiatement après une publication.
\end{itemize}

\subsection{Améliorations UX/UI et Accessibilité (Front-end)}
\begin{itemize}[label=\textcolor{primary}{\textbullet}]
    \item \textbf{Dates relatives :} Le \texttt{PhotoPostAdapter} calcule désormais des dates relatives (``Il y a 2h'', ``Hier'', ``Il y a 3j'') au lieu de dates absolues, conformément aux standards UX des réseaux sociaux.
    
    \item \textbf{Description dans les cartes :} Chaque carte de publication dans le fil affiche maintenant les 2 premières lignes de la description (tronquées par ellipsis), donnant plus de contexte sans alourdir l'interface.
    
    \item \textbf{Accessibilité TalkBack :} Ajout systématique des attributs \texttt{contentDescription} sur toutes les \texttt{ImageView} interactives (boutons retour, recherche, profil) et sur les éléments de la liste de publications. Les icônes décoratives sont marquées \texttt{importantForAccessibility="no"} pour éviter les doublons vocaux. La zone minimale de toucher des boutons a été portée à 48dp (recommandation Material Design).
    
    \item \textbf{Thème unifié :} Les deux thèmes (\texttt{Theme.Traveling} et \texttt{Theme.TravelPath}) utilisent maintenant les couleurs de charte définitives \texttt{\#7C4DFF}/\texttt{\#6C63FF} au lieu des couleurs Material par défaut. La barre de navigation système est également colorée en cohérence.
    
    \item \textbf{Placeholder élégant :} Un drawable vectoriel \texttt{placeholder\_post.xml} remplace le générique \texttt{ic\_menu\_gallery} pour les publications sans photo.
    
    \item \textbf{Animation d'entrée :} Les deux cartes de l'écran d'accueil (\texttt{HomeFragment}) se révèlent avec une animation \textit{fade + slide-up} décalée, donnant une impression de dynamisme à l'ouverture de l'application.
\end{itemize}

\section{Ce qu'il reste à faire / Pistes d'évolutions}

\begin{itemize}[label=\textcolor{primary}{\textbullet}]
    \item \textbf{Cloud Persistant et Synchronisation distante :}\\
    Actuellement, l'intégralité du réseau social TravelShare est restreint au périmètre confiné de la machine (SQLite). Pour valider le concept ``Communautaire'', il faudrait implémenter une API REST ou brancher un back-end as-a-Service, tel que \textit{Firebase (Firestore \& Storage)} ou  un back-end de type \textit{PostgreSQL/Spring Boot}, afin que le flux de publications et de profils se synchronise en direct entre deux téléphones distincts.

    \item \textbf{Optimisation du Routage (Temps-Réel) :}\\
    Le dessin de l'itinéraire sur la carte d'OSMDroid utilise majoritairement un calcul vectoriel pour tracer ses droits. L'utilisation et la substitution par une API dédiée et spécialisée (ex: \textit{GraphHopper Routing API} ou \textit{Google Directions API}) rendrait les horaires, les routes empruntées et les propositions de détours à pied beaucoup plus réalistes en se fiant aux topologies exactes des routes piétonnes.

    \item \textbf{Notifications Push :}\\
    Les interrupteurs de notification dans l'interface sont actuellement visuels uniquement. L'implémentation complète nécessiterait \textit{Firebase Cloud Messaging} ou un \texttt{WorkManager} Android pour déclencher des notifications système réelles lors de nouveaux likes, commentaires ou publications dans les groupes suivis.
\end{itemize}

\end{document}
